\chapter{Literature Integration and Evidence-Based Interpretation}

\WeekHeader{8}{Literature Integration and Evidence-Based Interpretation}{Connect database evidence to peer-reviewed context.}{Annotated evidence table and interpretation memo.}

\section{Why This Matters}
Database enrichment helps students find leads, but publication-quality interpretation requires literature context. A student must distinguish between a chemical being listed in a database, a chemical being measured in a system, and a chemical having a demonstrated role in a biological or environmental process.

\section{Evidence Ladder}
\begin{tabularx}{\textwidth}{p{0.28\textwidth}Y}
\toprule
\textbf{Evidence level} & \textbf{Interpretation strength}\\
\midrule
Tentative GC-MS match & Suggests possible presence; requires caution and preferably standards.\\
Database identifier or property & Supports identity, descriptors, and grouping.\\
Source annotation & Provides context from curated or aggregated source text.\\
Pathway, reaction, or occurrence evidence & Suggests biochemical, ecological, or taxonomic context.\\
Peer-reviewed measurement in similar system & Stronger support for relevance to the student's project.\\
Confirmed standard or targeted validation & Strongest support for chemical presence in the specific sample.\\
\bottomrule
\end{tabularx}

\section{Guided Lab}
Students select three high-priority findings from their Week 7 figures. For each finding, they fill an evidence table:
\begin{itemize}
  \item Chemical.
  \item Trait or pattern.
  \item uafR table and field.
  \item Source database.
  \item Literature citation.
  \item Interpretation.
  \item Limitation.
  \item Follow-up test.
\end{itemize}

\section{Finding Literature}
Use database records to find search terms, not conclusions. Good searches combine:
\begin{itemize}
  \item Chemical name or CID.
  \item System, organism, tissue, material, food, environment, disease, or exposure.
  \item Trait term, pathway, reaction, hazard code, sensory term, or measured property.
  \item Analytical method such as GC-MS, LC-MS, headspace, metabolomics, or targeted assay.
\end{itemize}

\section{Independent Extension}
Students write a two-page interpretation memo with:
\begin{itemize}
  \item One central result.
  \item Two supporting figures or tables.
  \item Three literature sources.
  \item One paragraph on uncertainty.
  \item One proposed follow-up experiment or analysis.
\end{itemize}

\section{Deliverables}
\begin{tabularx}{\textwidth}{p{0.25\textwidth}YY}
\DeliverableTableHeader
Annotated evidence table & Links database fields, literature, interpretation, and limits. & CSV or spreadsheet\\
Interpretation memo & Explains one result with evidence and caveats. & PDF\\
Reference library & Contains cleaned citations for final report. & BibTeX, Zotero, or plain bibliography\\
\end{tabularx}

\section{Quality Checklist}
\begin{checklistbox}{Before Week 9}
\begin{itemize}
  \item Claims are tied to a table, figure, or citation.
  \item The memo distinguishes tentative identity from confirmed presence.
  \item Literature sources are relevant to the system, not just the compound.
  \item Uncertainty is written clearly instead of hidden.
\end{itemize}
\end{checklistbox}
